% Quantum reservoir computing (TikZ).
% Single-column layout; classical (linear-regression) nodes use the Fig 1
% classical-NN blue. Body font via rmfamily.
\begingroup
% Palette from Fig 8: orange = inputs, green = reservoir, violet = readout/classical
\definecolor{inO}{RGB}{255,167,51}
\definecolor{resG}{RGB}{38,200,110}
\definecolor{outB}{RGB}{99,110,200}
\definecolor{mtrV}{RGB}{99,110,200}
\tikzset{
  innode/.style={fill=inO, draw=inO!60!black, line width=0.7pt},
  resnode/.style={fill=resG, draw=resG!60!black, line width=0.7pt},
  bluenode/.style={fill=outB, draw=outB!60!black, line width=0.7pt},
  sig/.style={black, line width=0.7pt, -{Stealth[length=2mm]}},
  res/.style={black, line width=0.7pt, -{Stealth[length=1.8mm]}},
  mtr/.style={draw=mtrV!60!black, line width=0.7pt, fill=mtrV, minimum width=0.7cm,
              minimum height=0.48cm, inner sep=1pt},
}
\begin{tikzpicture}[x=1cm,y=1cm, every node/.style={font=\rmfamily}]
% headers
\node at (0.55,3.3) {Inputs};
\node at (2.9,3.3)  {Reservoir};
\node at (5.75,3.3) {Readout};
\node[align=center] at (7.45,3.45) {Linear-\\Regression};
% reservoir internal edges (curved, drawn first)
\draw[res] (1.6,2.55) to[bend left=25] (2.5,2.25);
\draw[res] (2.5,2.25) to[bend left=20] (3.5,2.6);
\draw[res] (3.5,2.6) to[bend left=15] (4.4,2.25);
\draw[res] (1.95,1.5) to[bend right=20] (2.8,1.35);
\draw[res] (2.8,1.35) to[bend left=25] (2.5,2.25);
\draw[res] (2.8,1.35) to[bend right=15] (3.65,0.9);
\draw[res] (2.05,0.45) to[bend right=20] (2.7,-0.25);
\draw[res] (2.7,-0.25) to[bend right=15] (3.35,-0.6);
\draw[res] (3.35,-0.6) to[bend right=20] (4.05,0.2);
\draw[res] (3.65,0.9) to[bend left=15] (4.05,0.2);
\draw[res] (4.05,0.2) to[bend left=10] (4.6,1.0);
\draw[res] (1.95,1.5) to[bend left=25] (1.6,2.55);
\draw[res] (2.05,0.45) to[bend left=15] (1.95,1.5);
\draw[res] (1.7,-0.75) to[bend right=20] (2.7,-0.25);
\draw[res] (3.5,2.6) to[bend right=20] (2.8,1.35);
\draw[res] (4.4,2.25) to[bend left=20] (4.88,2.3);
\draw[res] (4.6,1.0) to[bend left=10] (4.88,1.28);
\draw[res] (4.05,0.2) to[bend right=10] (4.88,0.3);
\draw[res] (3.35,-0.6) to[bend right=15] (4.88,-0.7);
\draw[res] (4.4,2.25) to[bend right=25] (4.6,1.0);
% inputs -> reservoir
\foreach \y in {2.4,1.7,1.0,0.3,-0.4}{
  \node[innode, circle, minimum size=0.4cm] at (0.55,\y) {}; }
\draw[sig] (0.76,2.4) -- (1.4,2.5);
\draw[sig] (0.76,2.4) -- (1.75,1.6);
\draw[sig] (0.76,1.7) -- (1.42,2.42);
\draw[sig] (0.76,1.7) -- (1.75,1.52);
\draw[sig] (0.76,1.0) -- (1.75,1.42);
\draw[sig] (0.76,1.0) -- (1.85,0.5);
\draw[sig] (0.76,0.3) -- (1.85,0.42);
\draw[sig] (0.76,0.3) -- (1.5,-0.7);
\draw[sig] (0.76,-0.4) -- (1.5,-0.73);
\draw[sig] (0.76,-0.4) -- (1.85,0.36);
% reservoir nodes
\foreach \x/\y in {1.6/2.55, 2.5/2.25, 3.5/2.6, 1.95/1.5, 2.8/1.35,
                   2.05/0.45, 2.7/-0.25, 3.65/0.9, 3.35/-0.6, 4.4/2.25,
                   4.05/0.2, 1.7/-0.75, 4.6/1.0}{
  \node[resnode, circle, minimum size=0.4cm] at (\x,\y) {}; }
% read-out chain: output nodes -> meters -> blue nodes
\foreach \y in {2.3,1.3,0.3,-0.7}{
  \node[resnode, circle, minimum size=0.4cm] at (5.1,\y) {};
  \draw[sig] (5.31,\y) -- (5.57,\y);
  \node[mtr] at (5.97,\y) {};
  \draw[black, line width=0.6pt] (5.79,\y-0.1) arc[start angle=180, end angle=0, radius=0.18];
  \draw[black, line width=0.6pt] (5.97,\y-0.1) -- (6.09,\y+0.11);
  \draw[sig] (6.37,\y) -- (6.67,\y);
  \node[bluenode, circle, minimum size=0.4cm] at (6.9,\y) {};
  \node[bluenode, circle, minimum size=0.4cm] at (8.05,\y) {};
}
% full bipartite linear regression
\foreach \ya in {2.3,1.3,0.3,-0.7}{ \foreach \yb in {2.3,1.3,0.3,-0.7}{
  \draw[black, line width=0.6pt] (7.11,\ya) -- (7.84,\yb); }}
\end{tikzpicture}%
\endgroup
